
\documentclass{article}
\usepackage{colortbl}
\usepackage{makecell}
\usepackage{multirow}
\usepackage{supertabular}

\begin{document}

\newcounter{utterance}

\centering \large Interaction Transcript for game `guesswhat', experiment `Level\_3', episode 5 with reka{-}flash{-}3.1{-}t0.0.
\vspace{24pt}

{ \footnotesize  \setcounter{utterance}{1}
\setlength{\tabcolsep}{0pt}
\begin{supertabular}{c@{$\;$}|p{.15\linewidth}@{}p{.15\linewidth}p{.15\linewidth}p{.15\linewidth}p{.15\linewidth}p{.15\linewidth}}
    \# & $\;$A & \multicolumn{4}{c}{Game Master} & $\;\:$B\\
    \hline

    \theutterance \stepcounter{utterance}  
    & & \multicolumn{4}{p{0.6\linewidth}}{
        \cellcolor[rgb]{0.9,0.9,0.9}{
            \makecell[{{p{\linewidth}}}]{
                \texttt{\tiny{[P1$\langle$GM]}}
                \texttt{Let's play a guessing game. Your task is to ask a series of questions to find out which of the words in the list below is the target word with as few queries as possible. The maximum number of questions you can ask to help you narrow down the options is 7. Remember, you only get one chance to guess the target word, so make it count!} \\
\\ 
\texttt{The rules of the game are: } \\
\\ 
\texttt{1. Only ask one question per turn. DO NOT ASK MORE THAN ONE QUESTION.} \\
\\ 
\texttt{2. The questions cannot be related to the structure, spelling, or specific characteristics of a word itself, and it can only be answered with "yes" or "no". } \\
\\ 
\texttt{3. Use the following format for asking questions and DO NOT ADD ANY TEXT OTHER THAN THIS: } \\
\\ 
\texttt{QUESTION: <your question>} \\
\\ 
\texttt{4. You have only one opportunity to guess the target word. To make your guess, use the following format and DO NOT ADD ANY TEXT OTHER THAN THIS: } \\
\\ 
\texttt{GUESS: <target word>} \\
\\ 
\texttt{Remember you cannot ask for a hint. Make your question brief. Let’s start!} \\
\\ 
\texttt{List of candidate words: {[}'Turquoise', 'Opal', 'Topaz', 'Amethyst', 'Pearl', 'Ruby', 'Aquamarine', 'Emerald'{]}} \\
            }
        }
    }
    & \\ \\

    \theutterance \stepcounter{utterance}  
    & \multicolumn{4}{p{0.6\linewidth}}{
        \cellcolor[rgb]{0.9,0.9,0.9}{
            \makecell[{{p{\linewidth}}}]{
                \texttt{\tiny{[P1$\rangle$GM]}}
                \texttt{<reasoning>} \\
\texttt{Let's see. I need to figure out which gemstone is the target word from the list provided. The list is {[}'Turquoise', 'Opal', 'Topaz', 'Amethyst', 'Pearl', 'Ruby', 'Aquamarine', 'Emerald'{]}. The rules say I can ask up to 7 yes/no questions, but I have to make sure my final guess is correct. The questions can't be about the word's structure, spelling, or specific characteristics. So I can't ask about letters, syllables, or parts of the word. They have to be about the meaning or associations.} \\
\\ 
\texttt{First, I should think about common properties or associations of these gemstones. Maybe color? But some of them are different colors. For example, Ruby is red, Emerald is green, Amethyst is purple, etc. But the question can't be about color. Wait, the rules say the questions can't be related to the structure, spelling, or specific characteristics of the word itself. So maybe I can't ask about color either? Wait, the wording is a bit ambiguous. Let me check the rules again. Rule 2 says the questions cannot be related to the structure, spelling, or specific characteristics of a word itself. So "specific characteristics" might include color. Hmm. So maybe I can't ask about color. That complicates things.} \\
\\ 
\texttt{Alternatively, maybe I can ask about their properties or uses. For example, are they used in jewelry? Well, all of them are gemstones, so that's too broad. Maybe their hardness on the Mohs scale? But that's a specific characteristic. Wait, the rules say "specific characteristics of a word itself". Maybe hardness is a property of the gemstone, not the word. So maybe that's allowed. Let me think. The question has to be answerable with yes or no, and can't be about the word's structure or spelling. So if I ask, "Is the target word a precious gemstone?" the answer would be yes, but that's not helpful because all the words in the list are gemstones. So that's not useful.} \\
\\ 
\texttt{Alternatively, maybe think about their origin. For example, Turquoise comes from the Persian word. But how would that help? Maybe ask if the gemstone is named after a place. But that's a stretch. Alternatively, think about their colors again. Wait, maybe the question can be about the color,} \\
            }
        }
    }
    & & \\ \\

    \theutterance \stepcounter{utterance}  
    & & & \multicolumn{2}{p{0.3\linewidth}}{
        \cellcolor[rgb]{0.9,0.9,0.9}{
            \makecell[{{p{\linewidth}}}]{
                \texttt{\tiny{[GM$|$GM]}}
                \texttt{Invalid format. Guesser must use the form 'QUESTION: ' or 'GUESS: '.} \\
            }
        }
    }
    & & \\ \\

    \theutterance \stepcounter{utterance}  
    & & & \multicolumn{2}{p{0.3\linewidth}}{
        \cellcolor[rgb]{0.9,0.9,0.9}{
            \makecell[{{p{\linewidth}}}]{
                \texttt{\tiny{[GM$|$GM]}}
                \texttt{abort game} \\
            }
        }
    }
    & & \\ \\

\end{supertabular}
}

\end{document}
